% šablona na pro bakalářskou práci na FAI UTB
% šablona je úpravou původní šablony pro FAI UTB
% už neznám púvodního autora šablona už není na internetu
% upravil Martin Zajíc
% šablona je upravena aby šla bez problému přeložit pomocí TeXlive a překladače XeLaTeX
% 
\documentclass[a4paper,12pt]{article}
\usepackage{xltxtra}
\usepackage[czech]{babel}
\usepackage{indentfirst}
\usepackage{color}
\usepackage{colortbl}
\usepackage{array}
\usepackage{graphicx}
\usepackage{amsmath}
\usepackage{fancyhdr}
\usepackage{hyperref}
\usepackage{multirow}
\usepackage{tocloft}
\usepackage{textcase}
\usepackage{ifthen}
\usepackage{setspace}
\usepackage{ccaption}
\usepackage{sectsty}
\usepackage[srcstyle=leftnumhang,linenumbersep={\ }]{examplep}
% ************************ NASTAVENÍ AUTORA A NÁZVU DOKUMENTU ************************
\newcommand{\rok}{2011}
\newcommand{\jmeno}{Martin Zajíc}
\newcommand{\typprace}{Bakalářská práce}
\newcommand{\predmet}{Tvorba distribuce OS Linux}
\newcommand{\predmeten}{English}
% ************************ NASTAVENÍ TISKOVÉHO ZRCADLA ************************
\textheight=248mm
\textwidth=155mm
\voffset=-1.61cm
\oddsidemargin=0.96cm
\evensidemargin=0.96cm

% nastavení záhlaví
\headheight=0.5cm
\headsep=1cm

% nastavení zápatí
\footskip=1ex
\rhead{\thepage}
\cfoot{}
% "vypnout" poznámky na okrajích
\marginparpush=0mm
\marginparwidth=0mm
\marginparsep=0mm

\pagestyle{fancy}
% ****************** NASTAVENÍ HYPERPREF (BAREV ODKAZŮ, PDF DOKUMENTU) **************
\hypersetup{
    pdftoolbar=true,        % show Acrobat’s toolbar?
    pdfmenubar=true,        % show Acrobat’s menu?
    pdftitle={\typprace~-~\predmet },    % title
    pdfauthor={\jmeno},     % author
    pdfsubject={\predmet},   % subject of the document
    pdfcreator={\jmeno},   % creator of the document
    colorlinks=true,       % false: boxed links; true: colored links
    linkcolor=black,          % color of internal links
    citecolor=black,        % color of links to bibliography
    filecolor=black,      % color of file links
    urlcolor=black           % color of external links
}

% ****************** NASTAVENÍ PÍSMA, ODSTAVCE, ROVNIC, POZNÁMEK **************
\parindent=0em
\def\thefootnote{\arabic{footnote})}	% poznámka pod čarou se závorkou
\onehalfspacing % nastavím řádkování tímto způsobem nebo \renewcommand{\baselinestretch}{1.5} ??
\setlength{\parskip}{3pt}		% vertikální mezera mezi nadpisy
% *************************** NASTAVENÍ ČÍTAČŮ	 ********************************
\setcounter{tocdepth}{3} % do obsahu se ukládají pouze první dvě úrovně kapitol

% *************************** UŽIVATELSKÉ STYLY *******************************
\newcommand{\nn}[1]{\clearpage\section*{\texorpdfstring{\uppercase{#1}}{#1}}\addcontentsline{toc}{section}{\uppercase{#1}}}% styl nn = nečíslovaný nadpis (je vysázený v obsahu)
\newcommand{\nm}[1]{\clearpage\section*{\uppercase{#1}}}	% definujeme styl nm = nečíslovaný nadpis (není vysázený v obsahu)
\newcommand{\nmm}[1]{\section*{#1}} % definujeme styl nmm = nečíslovaný nadpis (není vysázený v obsahu) malými písmeny bez clearpage
\newcommand{\nns}[1]{\section*{\uppercase{#1}}}		% definujeme styl ns = nečíslovaný nadpis na stejné stránce (není vysázený v obsahu)
\newcommand{\sectionV}[1]{\section{\uppercase{#1}}}	%nadpis1 (section) velkým

\newcommand{\obr}[3]{% styl \obr pro obrázky
	\begin{figure}[h]
	\center\includegraphics[scale=#2]{#1}
	\caption{#3}
	\end{figure}
	}

\newcommand{\tab}[3]{% styl \tab pro tabulky
	\begin{table}[h]
	\caption{#1}
	\begin{center}
	\begin{tabular}{#2}
	#3
	\end{tabular}
	\end{center}
	\end{table}
	}

\newcommand{\tabpri}[3]{% styl \tabpri pro tabulky v příloze
	\begin{table}[h]
	\begin{center}
	#1
	\end{center}
	\begin{center}
	\begin{tabular}{#2}
	#3
	\end{tabular}
	\end{center}
	\end{table}
	}

\newcommand{\rov}[2][chybějici rovnice]{% styl \rov pro rovnice
	\begin{equation}
	#2
	\label{#1}
	\end{equation}
	}
	
\newcommand{\seznamobr}{% příkaz \seznamobr pro vysázení seznamu obrázků
	\addcontentsline{toc}{section}{\listfigurename}
	\clearpage
	\listoffigures
	\clearpage
	}

\newcommand{\seznamtab}{% příkaz \seznamtab pro vysázení seznamu obrázků
	\addcontentsline{toc}{section}{\listtablename}
	\clearpage
	\listoftables
	\clearpage
	}
	
\newcommand{\seznamlit}[1]{% příkaz \seznamlit pro vysázení seznamu literatury
	\addcontentsline{toc}{section}{\refname}
	\begin{thebibliography}{10}
	#1
	\end{thebibliography}}
	
\newcommand{\seznamzkr}{% příkaz \seznamzkr pro přípravu seznamu použitých zkratek a symbolů
	\nn{Seznam použitých symbolů a zkratek}
	}
	
\newcommand{\obsah}{% \obsah vysází obsah v daném místě
	\clearpage
	\thispagestyle{empty}
	\tableofcontents
	\clearpage
	\pagestyle{fancy}
	}

\renewcommand{\b}[1]{\textbf{#1}} % \b = tučně
\newcommand{\bi}[1]{\textbf{\textit{#1}}}	% \bi = tučná kurzíva
\renewcommand{\it}[1]{\textit{#1}}		% \it = kurzíva
	
% ******* NASTAVENÍ ZOBRAZENÍ PŘÍLOH -- SEZNAM, ČÍSLOVÁNÍ, VLASTNÍ STYL *******
\makeatletter % tímto příkazem dávám najevo, že budu editovat přímo příkazy ze šablony

% definice seznamu příloh - příkaz \listofappendices
\newcommand{\listofappendices}{%
	\newpage
	\setcounter{section}{0}
	\nn{Seznam příloh}
	\@restonecolfalse\if@twocolumn\@restonecoltrue\onecolumn\fi
	\@mkboth{LIST OF APPENDICES}{LIST OF APPENDICES}
	\@starttoc{loa}\if@restonecol\twocolumn\fi
	\pagestyle{empty}
	\thispagestyle{fancy}
	}

\def\ext@appendix{loa}
\def\tocname{loa}

% definice příkazu \priloha{nazev prilohy} pro vložení nové přílohy
\newcommand{\priloha}[1]{
	\clearpage
	\refstepcounter{section}
	\addtocontents{loa}{\protect\makebox[1.5cm][l]{P \@Roman\c@section.} #1\newline}
{\bf PŘÍLOHA P \@Roman\c@section. \uppercase{#1}}\par}

% v obsahu nastavím VELKÉ PÍSMENA pro styl část
\renewcommand{\part}[1]{
	\refstepcounter{part}
	\addcontentsline{toc}{section}{\thepart~~\uppercase{#1}}%
\clearpage
\normalfont
	\vspace*{9cm}
	\begin{center}\huge \bfseries\thepart. \uppercase{#1}\end{center}%
	\markboth{}{}\par
\nobreak
\clearpage	
}

% *************** NASTAVENÍ FORMÁTU ČÍSLOVÁNÍ OBRÁZKŮ A TABULEK ***************
\def\thefigure{\arabic{figure}}      % číslování obrázků typu (y)
\def\thetable{\arabic{table}}        % číslování tabulek typu (y)
\captiondelim{. } % změníme dvoutečku za Obr/Tab za tečku

% ***************************** VYNULOVÁNÍ ČÍTAČŮ *****************************
\@addtoreset{table}{section}    % vynulování čítače
\@addtoreset{figure}{section}   % vynulování čítače
\@addtoreset{footnote}{section} % vynulování čítače
	
\makeatother % a to je ukončení \makeatletter

% ************* ÚPRAVA VZHLEDU OBSAHU, SEZNAMU OBRÁZKÙ A TABULEK **************
% nastavení vertikální mezery před stylem část, nadpis 1--3
\setlength{\cftbeforepartskip}{3pt}
\setlength{\cftbeforesecskip}{3pt}
\setlength{\cftbeforesubsecskip}{3pt}
\setlength{\cftbeforesubsubsecskip}{0cm}
% odsazení zleva pro styl část, nadpis 1--3
\setlength{\cftpartindent}{0cm}
\setlength{\cftsecindent}{0cm}
\setlength{\cftsubsecindent}{0cm}
\setlength{\cftsubsubsecindent}{0cm}
% nastavení fontu pro styl část, nadpis 1--3
\renewcommand{\cftpartfont}{\small\bfseries}
\renewcommand{\cftsecfont}{\small\bfseries}
\renewcommand{\cftsubsecfont}{\scshape}
\renewcommand{\cftsubsubsecfont}{}
% odsazení čísla a textu titulku pro styl část, nadpis 1--3
\cftsetindents{part}{0cm}{1cm}
\cftsetindents{sec}{0cm}{1cm}
\cftsetindents{subsec}{1cm}{1.5cm}
\cftsetindents{subsubsec}{1.2cm}{1.7cm}
\cftsetindents{fig}{0cm}{1.5cm}
\cftsetindents{tab}{0cm}{1.5cm}
% nastavení vodící čáry pro styl část, nadpis 1--3, obrázky a tabulky
\renewcommand{\cftdot}{\ensuremath{.}} % tímto příkazem lze změnit vodící tečky v obsahu na jiný znak
\renewcommand{\cftpartleader}{\cftdotfill{0.3}}
\renewcommand{\cftsecleader}{\cftdotfill{0.3}}
\renewcommand{\cftsubsecleader}{\cftdotfill{0.3}}
\renewcommand{\cftsubsubsecleader}{\cftdotfill{0.3}}
\renewcommand{\cftfigleader}{\cftdotfill{0.3}}
\renewcommand{\cfttableader}{\cftdotfill{0.3}}
% změna fontu pro text "Obsah", "Seznam obrázků" a "Seznam tabulek"
\renewcommand{\cfttoctitlefont}{\normalsize\bfseries\thispagestyle{empty}}
\renewcommand{\cftloftitlefont}{\normalsize\bfseries\thispagestyle{fancy}}
\renewcommand{\cftlottitlefont}{\normalsize\bfseries\thispagestyle{fancy}}
\renewcommand{\cftfigpresnum}{Obr. }
\renewcommand{\cftfigaftersnum}{.}
\renewcommand{\cfttabpresnum}{Tab. }
\renewcommand{\cfttabaftersnum}{.}

% ************************ NASTAVENÍ FONTU PRO NADPISY ************************
\sectionfont{\normalsize}
\subsectionfont{\normalsize\bfseries} 
\subsubsectionfont{\normalsize\bfseries}
\paragraphfont{\small\bf}
% definice nového stylu \comment -- komentář k šabloně
\newcommand{\comment}[1]{\color{red}#1\color{black}}

% ****************************** KONEC PREAMBULE ******************************
\begin{document}
% ****************************** Vlastní definice ******************************
% ********* musí být na začátku dokumnetu v preambuli nefungují ****************
\renewcommand{\refname}{\texorpdfstring{\uppercase{Seznam použité literatury}}{Seznam použité literatury}}
\renewcommand*{\listfigurename}{\texorpdfstring{\uppercase{Seznam obrázků}}{Seznam obrázků}}
\renewcommand{\contentsname}{\texorpdfstring{\uppercase{Obsah}}{Obsah}}
\renewcommand*{\listtablename}{\texorpdfstring{\uppercase{Seznam tabulek}}{Seznam tabulek}}
\renewcommand*{\tablename}{Tab.}
\renewcommand\figurename{Obr.}
%********************************* Hlavička ************************************
\lhead{\b{UTB ve Zlíně, Fakulta aplikované informatiky}}

\thispagestyle{empty}
\voffset=-2.01cm\evensidemargin=0pt\oddsidemargin=0cm\parindent=0pt\headsep=0pt\headheight=0pt\parskip=0pt\textheight=272mm\textwidth=200mm
\renewcommand{\baselinestretch}{0}
%********************************* Úvodní stránka s logem fakulty ********************
\setlength{\unitlength}{1mm}
\begin{picture}(-10,8)
\put(0,-50){\makebox(170,50){}}
%*******************************Uvodni stranka z nazvem prace na dva radky*********
%\put(0,-100){\makebox(170,50){\fontsize{24}{1}\b{\textsf{Zálohovací server založený na}}}}
%\put(0,-110){\makebox(170,50){\fontsize{24}{1}\b{\textsf{GNU/Linux Debian}}}}
%\put(0,-130){\makebox(170,50){\fontsize{20}{1}\b{\textsf{Backup server based on GNU/Linux Debian}}}}
%\put(0,-150){\makebox(170,25){\fontsize{20}{1}\b{\textsf{\jmeno}}}}
%*******************************konec uvodni stranka na dva radky***********************
\put(0,-115){\makebox(170,50){\fontsize{24}{1}\b{\textsf{\predmet}}}}
\put(0,-130){\makebox(170,50){\fontsize{22}{1}\b{\textsf{\predmeten}}}}
\put(0,-150){\makebox(170,25){\fontsize{20}{1}\b{\textsf{\jmeno}}}}
\put(82.2,-223.3){\makebox(84,16.4){\includegraphics[width=84mm]{fai_logo_cz.png}}}
\put(0,-205){\linethickness{1pt}\line(1,0){170}}
\put(4,-215){\makebox(69.5,4.5)[l]{\noindent\fontsize{16}{1}\b{\textsf{\typprace}}}}
\put(4,-220){\makebox(69.5,4.5)[l]{\noindent\fontsize{16}{1}\b{\textsf{\rok}}}}
\put(0,-225){\linethickness{1pt}\line(1,0){170}}
\put(75,-223.3){\linethickness{1pt}\line(0,1){16.4}}
\end{picture}
%********************************* nascanované zadání ************************************
\clearpage
\thispagestyle{empty}
\voffset=-1.61cm\evensidemargin=0.96cm\oddsidemargin=0.96cm\headsep=1cm\headheight=0.5cm\setlength{\parskip}{3pt}\textheight=248mm\textwidth=155mm
***nascannované zadání s. 1***
\obr{fai_logo_cz.png}{1}{logo}
\clearpage
\thispagestyle{empty}
***nascannované zadání s. 2***
\obr{fai_logo_cz.png}{1}{logo}
%********************************* Abstrakt ************************************
\clearpage
\thispagestyle{empty}
\nm{\uppercase{Abstrakt}}
\vspace{1cm}

\emph{Klíčová slova:}

\vspace{3cm}

\nns{\uppercase{Abstract}}

\vspace{1cm}

\emph{Keywords:}
%********************************* poděkování, motto, úryvky knih, básní atp. ************************************
\clearpage
\thispagestyle{empty}
poděkování, motto, úryvky knih, básní atp.
%********************************* Prohlášení ************************************
\clearpage
\thispagestyle{empty}
\clearpage
\nmm{Prohlašuji, že}
\begin{itemize}\fontsize{10}{1}
\normalsize
\item beru na vědomí, že odevzdáním bakalářské práce souhlasím se zveřejněním své práce podle zákona č. 111/1998 Sb. o vysokých školách a o změně a doplnění dalších zákonů (zákon o vysokých školách), ve znění pozdějších právních předpisů, bez ohledu na výsledek obhajoby;
\item beru na vědomí, že bakalářská práce bude uložena v elektronické podobě v univerzitním informačním systému dostupná k prezenčnímu nahlédnutí, že jeden výtisk bakalářské práce bude uložen v příruční knihovně Fakulty aplikované informatiky Univerzity Tomáše Bati ve Zlíně a jeden výtisk bude uložen u vedoucího práce; 
\item byl/a jsem seznámen/a s tím, že na moji bakalářskou práci se plně vztahuje zákon č. 121/2000 Sb. o právu autorském, o právech souvisejících s právem autorským a o změně některých zákonů (autorský zákon) ve znění pozdějších právních předpisů, zejm. § 35 odst. 3;
\item beru na vědomí, že podle § 60 odst. 1 autorského zákona má UTB ve Zlíně právo na uzavření licenční smlouvy o užití školního díla v rozsahu § 12 odst. 4 autorského zákona;
\item beru na vědomí, že podle § 60 odst. 2 a 3 autorského zákona mohu užít své dílo –bakalářskou práci nebo poskytnout licenci k jejímu využití jen s předchozím písemným souhlasem Univerzity Tomáše Bati ve Zlíně, která je oprávněna v takovém případě ode mne požadovat přiměřený příspěvek na úhradu nákladů, které byly Univerzitou Tomáše Bati ve Zlíně na vytvoření díla vynaloženy (až do jejich skutečné výše);
\item beru na vědomí, že pokud bylo k vypracování bakalářské práce
využito softwaru poskytnutého Univerzitou Tomáše Bati ve Zlíně nebo jinými
subjekty pouze ke studijním a výzkumným účelům (tedy pouze k nekomerčnímu
využití), nelze výsledky bakalářské práce využít ke komerčním
účelům;
\item beru na vědomí, že pokud je výstupem bakalářské práce jakýkoliv softwarový produkt, považují se za součást práce rovněž i zdrojové kódy, popř. soubory, ze kterých se projekt skládá. Neodevzdání této součásti může být důvodem k neobhájení práce.
\end{itemize}
\nmm{Prohlašuji,}
\begin{itemize}
 \normalsize
\item že jsem na bakalářské práci pracoval samostatně a použitou literaturu jsem citoval. V případě publikace výsledků budu uveden jako spoluautor.
\item že odevzdaná verze bakalářské práce a verze elektronická nahraná do IS/STAG jsou totožné.
\end{itemize}
\vspace{1cm}
\begin{tabular}{lc}
Ve Zlíně &\hspace{9cm} .......................\\
~ &\hspace{9cm} podpis diplomanta
\end{tabular}

% obrázky \obr{soubor.png}{scale}{popisek}
% tabulky \tab{popisek}{nastaveni}{obsah}
% tabulky v příloze \tabpri{vrch tabulky}{nastaveni}{obsah}
% rovnice \rov{popisek}{obsah}
% tučně \b{tučné}
% tučná kurzíva \bi{tučná kurzíva}
% kurzíva \it{kurzíva}
%********************************* Generování obsahu ************************************
\obsah
%********************************* úvod ************************************
\nn{Úvod}
text

%********************************* Vlastní obsah ************************************
%********************************* Teoretická část ************************************
\part{Teoretická část}

\sectionV{Nadpis}
text
\sectionV{Nadpis}
\subsection{Podnadpis}
text~\cite{chmel}

\subsubsection{Podpodnadpis}
text

\subsubsection{Podpodnadpis}
text
%********************************* Analytická část ************************************
\part{Analytická část}
\sectionV{Nadpis}
\subsection{Podnadpis}
%********************************* Projektová část ************************************
\part{Projektová část}

\sectionV{Nadpis}
\subsection{Podnadpis}

%********************************* Závěr ************************************
\nn{Závěr}
text
%********************************* Seznam Literatury ************************************
\seznamlit{
\bibitem{chmel}{\it{Český chmel: atlas odrůd}} [online]. [cit. 2004-10-15]. Dostupný z~WWW: \href{http://www.beer.cz/humulus/}{http://www.beer.cz/humulus/}.
\bibitem{chmel1}{\it{Český chmel: atlas odrůd}} [online]. [cit. 2004-10-15]. Dostupný z~WWW: \href{http://www.beer.cz/humulus/}{http://www.beer.cz/humulus/}.
\bibitem{chmel2}{\it{Český chmel: atlas odrůd}} [online]. [cit. 2004-10-15]. Dostupný z~WWW: \href{http://www.beer.cz/humulus/}{http://www.beer.cz/humulus/}.
}
%********************************* Seznam zkratek ************************************
\seznamzkr

\begin{tabular}{ll}
ABC & Význam zkratky \\
\end{tabular}
%********************************* Seznam obrázků ************************************
\seznamobr
%********************************* Seznam tabulek ************************************
\seznamtab
%********************************* Seznam příloh ************************************
\listofappendices

\priloha{Název přílohy}
%\section*{PŘÍLOHA P I: ŠABLONA SLOVNĚ.}

\end{document}
